\documentclass[11pt]{article}

% ============================================================
% Packages
% ============================================================
\usepackage[utf8]{inputenc}
\usepackage[margin=1in]{geometry}
\usepackage{amsmath,amssymb,amsthm}
\usepackage{graphicx}
\usepackage{booktabs}
\usepackage{xcolor}
\usepackage{hyperref}
\usepackage{tikz}
\usepackage{enumitem}
\usepackage{pgfplots}

\pgfplotsset{compat=1.18}

% ============================================================
% Document Formatting
% ============================================================
\setlength{\parskip}{0.6em}
\setlength{\parindent}{0pt}

% ============================================================
% List Formatting
% ============================================================
\setlist[enumerate]{
    leftmargin=2em,
    itemsep=0.25em
}

% ============================================================
% Solution Environment
% ============================================================
\newenvironment{solution}
    {\par\noindent\textbf{Solution.}\par}
    {\par}

% ============================================================
% Link Styling
% ============================================================
\hypersetup{
    colorlinks=true,
    linkcolor=blue,
    urlcolor=cyan
}

% ============================================================
% Metadata
% ============================================================
\title{Ordinary Differential Equations: Abridged Notes}
\author{Chris Ritter}
\date{September 22, 2026}

% ============================================================
% Document
% ============================================================

\begin{document}
\maketitle
\begin{center}
\small\emph{These abridged notes were reconstructed afterward because I did not take
direct notes during class. They are not a record of the actual lecture.}
\end{center}
\tableofcontents
\newpage

\section{Topic Summary}
Two solutions of a second-order homogeneous linear equation form a
fundamental pair when they are linearly independent. Their Wronskian
provides a test for independence. Abel's formula predicts how the
Wronskian changes without requiring explicit solutions. These facts
also summarize the existence and uniqueness ideas useful for review:
a second-order linear initial-value problem has a unique solution on
an interval where its normalized coefficients are continuous.

\section{The Wronskian and Abel's Formula}
For differentiable functions \(y_1,y_2\), define
\[
W(y_1,y_2)(t)=
\begin{vmatrix}y_1&y_2\\y_1'&y_2'\end{vmatrix}
=y_1y_2'-y_1'y_2.
\]
For two solutions of \(y''+p(t)y'+q(t)y=0\), Abel's formula says
\[
\boxed{W(t)=W(t_0)\exp\left(-\int_{t_0}^{t}p(s)\,ds\right).}
\]
Thus on an interval where \(p,q\) are continuous, \(W\) is either
nowhere zero or identically zero. A nonzero Wronskian means the pair
is independent and spans all solutions of the homogeneous equation.
When the leading coefficient vanishes, normalize only on intervals
that do not cross that point.

\subsection{Example: Compute a Wronskian Directly}
\textbf{Adapted from Boyce et al., 11th ed., \S3.2, Problem 3.}
Find the Wronskian of \(y_1=e^{-2t}\) and \(y_2=te^{-2t}\).
\begin{solution}
Since \(y_1'=-2e^{-2t}\) and
\(y_2'=(1-2t)e^{-2t}\),
\[
\begin{aligned}
W&=e^{-2t}(1-2t)e^{-2t}
-(-2e^{-2t})(te^{-2t})\\
&=\boxed{e^{-4t}}.
\end{aligned}
\]
It never vanishes, so the two functions are linearly independent on
all of \(\mathbb R\).
\end{solution}

\subsection{Example: Determine a Wronskian Without Solving}
\textbf{Adapted from Boyce et al., 11th ed., \S3.2, Problem 25.}
For two solutions of Legendre's equation
\[
(1-x^2)y''-2xy'+\alpha(\alpha+1)y=0,
\]
find the possible form of their Wronskian.
\begin{solution}
Away from \(x=\pm1\), divide by \(1-x^2\). The coefficient of
\(y'\) is \(p(x)=-2x/(1-x^2)\). Abel's formula gives
\[
W(x)=C\exp\left(\int\frac{2x}{1-x^2}\,dx\right)
=C\exp\bigl(-\ln|1-x^2|\bigr).
\]
On any interval excluding \(-1\) and 1, absorb the fixed sign of
\(1-x^2\) into the constant to obtain
\[
\boxed{W(x)=\frac{C}{1-x^2}.}
\]
If the solutions form a fundamental pair then \(C\ne0\). The three
allowed maximal coefficient intervals are
\(( -\infty,-1)\), \((-1,1)\), and \((1,\infty)\).
\end{solution}
\end{document}
